\uuid{5eFT}
\exo7id{1623}
\titre{exo7 1623}
\auteur{barraud}
\organisation{exo7}
\datecreate{2003-09-01}
\isIndication{false}
\isCorrection{false}
\chapitre{Réduction d'endomorphisme, polynôme annulateur}
\sousChapitre{Diagonalisation}
\cpge{mpsi-mp2i.determinants.d, pcsi.calcul-matriciel-systemes-lineaires.d, ptsi.calcul-matriciel-systemes-lineaires.d, bcpst1.matrices, mp-mpi.reduction-endomorphismes.c, pc.reduction.b, psi.reduction-endomorphismes.b, pt.reduction-endomorphismes.b, mp-mpi.reduction-endomorphismes.d, pc.reduction.c, psi.reduction-endomorphismes.c, pt.reduction-endomorphismes.c, bcpst2.valeurs-propres-vecteurs-propres.a, bcpst2.valeurs-propres-vecteurs-propres.b}
\module{Algèbre}
\niveau{L2}
\difficulte{}

\contenu{
\question{
Pour quelles valeurs de $(a,b,c)\in\C^{2}$ la matrice $A=
  \begin{pmatrix}
    1&a&1&0\\
    0&1&b&2\\
    0&0&2&c\\
    0&0&0&2
  \end{pmatrix} 
$ est-elle diagonalisable~? On ne cherchera pas à réduire explicitement
$A$.
}
}
